\documentclass{article}
\usepackage[UTF8]{ctex}
\usepackage{amsmath,amssymb,amsfonts}
\usepackage{enumitem}
\usepackage{geometry}
\geometry{a4paper, margin=1in}
% === 错题本排版优化（不参与题目标识） ===
\setlength{\parindent}{0pt}
\setlength{\parskip}{0.35em plus 0.10em minus 0.08em}
\setlength{\emergencystretch}{3em}
\setlist{itemsep=0.25em, topsep=0.40em, parsep=0pt}
\allowdisplaybreaks[2]
% === 错题本排版优化结束 ===
\begin{document}
\section*{华南理工大学2026年数学分析真题}
\section*{157．华南理工大学2026年数学分析试卷}

\subsection{题目1}
1．设 \(x_{0}=2, x_{n+1}=\frac{1}{2}\left(x_{n}+\frac{3}{x_{n}}\right)(n=0,1, \cdots)\) ，证明： \(\lim _{n \rightarrow \infty} x_{n}=\sqrt{3}\) ．

\subsection{题目2}
2．\(f(x)=x^{2} \sin x\) 在 \(0 \leq x<+\infty\) 上是否一致连续？给出证明．

\begin{itemize}

\subsection{题目3}
\end{itemize}

\subsection{题目3}
原题边界“3．”：\(f\) 在 \([0,2]\) 上二阶连续可微，且 \(\left|f^{\prime \prime}(x)\right| \leq M\) 。若 \(f\) 在 \((0,2)\) 内能取到最大值，证明：
\[
\left|f^{\prime}(0)+f^{\prime}(2)\right| \leq 2 M
\]

\subsection{题目4}

\subsection{题目4}
原题边界“4．”：已知 \(f(x)\) 在 \([A, B]\) 上黎曼可积，证明： \(\lim _{h \rightarrow 0} \int_{a}^{b}|f(x+h)-f(x)| \mathrm{d} x=0\) ，其中 \(A<a<b<B\) ．

\subsection{题目5}
5．设 \(f(t)\) 在 \((-\infty,+\infty)\) 上连续且绝对可积．证明：\(g(w)=\int_{-\infty}^{+\infty} f(t) \cos (w t) \mathrm{d} t\) 在 \(w \in(-\infty,+\infty)\)上有界且一致连续。

\subsection{题目6}
6．判断 \(\sum_{n=2}^{\infty} \ln \left(1+\frac{(-1)^{n}}{n^{\alpha}}\right)(\alpha>0)\) 的绝对收敛性，条件收敛性，发散性。

\subsection{题目7}
7．解答如下问题：

\begin{itemize}
  \item[（1）] 求 \(|\sin x|\) 在 \((-\infty,+\infty)\) 上的傅里叶级数，该级数是否一致收敛？
  \item[（2）] 若 \(f(x)\) 在 \([a, b]\) 上黎曼可积，证明： \(\lim _{\lambda \rightarrow+\infty} \int_{a}^{b} f(x)|\sin \lambda x| \mathrm{d} x=\frac{2}{\pi} \int_{a}^{b} f(x) \mathrm{d} x\) ．
\end{itemize}

\subsection{题目8}
8．记区域 \(D: 0 \leq x \leq 1,0 \leq y \leq 1\) ，证明：\(e^{x y}+\ln (1+x y) \geq 1+x y\) ，并说明取等条件．

\subsection{题目9}
9．已知

\[
f(x, y)= \begin{cases}\frac{x^{3}}{x^{2}+y^{2}}, & x^{2}+y^{2} \neq 0 \\ 0, & x^{2}+y^{2} \neq 0\left(\text { 应该是 } x^{2}+y^{2}=0\right)\end{cases}
\]

（1）若 \(x=x(t), y=y(t)\) 是过原点的任意可微曲线（即 \(x^{2}(0)+y^{2}(0)=0\) ，且 \(\left(x^{\prime}(t)\right)^{2}+\left(y^{\prime}(t)\right)^{2}>0\) ），证明：\(f(x(t), y(t))\) 可微．\\
（2）求证 \(f(x, y)\) 在 \((0,0)\) 处不可微．

\subsection{题目10}
1 ．求 \(\lim _{x \rightarrow 0} \frac{(5+x)^{x}-5^{x}}{x^{2}}\) ．

\subsection{题目11}
2 ．求

\[
I=\int_{L}\left(x^{2}+z^{2}\right) \mathrm{d} y+\left(y^{2}+x^{2}\right) \mathrm{d} z
\]

\(L\) 是 \(x^{2}+y^{2}+z^{2}=2 a x\) 与 \(y^{2}+x^{2}=2 b x\) 的交线，方向是经由 \(y>0\) 的方向再回到原点（ \(a>b\) ）。

\subsection{题目12}
3．已知 \(f(x, y, z)=x^{2}+y^{2}-z, \sum\) 是一光滑封闭曲面，\(\left(x_{0}, y_{0}, z_{0}\right)\) 是单位球体上的一点， \(\mathbf{n}\) 是 \(\sum\) 的单位外法向量，计算

\[
I=\iint_{\Sigma}\left(f \frac{\cos (r, n)}{r^{2}}+\frac{1}{r} \frac{\partial f}{\partial n}\right) \mathrm{d} S
\]

其中 \(r=\left(x-x_{0}, y-y_{0}, z-z_{0}\right), r=|\mathbf{r}|\) 。（ \(I\) 应该无误，其余 \(\sum, \mathbf{n}\) 等有关叙述可能不太准则）
\end{document}
